% Compile with: uplatex file.tex && dvipdfmx file.dvi
%           or: latexmk -uplatex file.tex
%\documentclass[uplatex,dvipdfmx]{jsarticle}
\documentclass[dvipdfmx,a4paper,11pt]{jsarticle}
\usepackage{amsmath, amssymb}
\usepackage{tikz}

\title{scenario\_bcs\_gap\_equation}
\date{\today}

\begin{document}
\maketitle

BCSハミルトニアンを組み立てるための演算子・パラメータを準備する。$c$ は運動量 $k$・スピン $\sigma$（$u=\uparrow$, $d=\downarrow$）を持つ電子の消滅演算子、$g$ は Bogoliubov 準粒子の消滅演算子。$u, v$ は変換係数、$\xi$ は運動エネルギー、$\Delta$ は超伝導ギャップ。

c を電子の消滅演算子として定義して

\[
  c
\]

g をフェルミオン演算子として定義して

\[
  g
\]

u, v をスカラーパラメータとして定義して

\[
  u, v
\]

xi, Delta をスカラーパラメータとして定義して

\[
  \xi, \Delta
\]

平均場BCSハミルトニアンを登録する。

$H_BCS = \xi (c'(k,u) c(k,u) + c'(-k,d) c(-k,d)) + \Delta (c'(k,u) c'(-k,d) + c(-k,d) c(k,u))$

\begin{equation}
  H_BCS = \xi (c^{\dagger}(k,u) c(k,u) + c^{\dagger}(-k,d) c(-k,d)) + \Delta (c^{\dagger}(k,u) c^{\dagger}(-k,d) + c(-k,d) c(k,u))
  \label{eq:E1}
\end{equation}

運動方程式で $c(k,u)$ が H\_BCS とどう結びついているかを確認する。

c(k,u) の Heisenberg 方程式 E1

\begin{equation}
  \frac{d c(k,u)}{d t} = - \Delta \mathit{i} {c}_{-k,d}^{\dagger} - \xi \mathit{i} {c}_{ku}
  \label{eq:E2}
\end{equation}

Bogoliubov 変換を Nambu スピノル形式の4本まとめて定義する。

$g(k,u) = u c(k,u) - v c'(-k,d)$

\begin{equation}
  g(k,u) = u c(k,u) - v c^{\dagger}(-k,d)
  \label{eq:E3}
\end{equation}

$g(-k,d) = u c(-k,d) + v c'(k,u)$

\begin{equation}
  g(-k,d) = u c(-k,d) + v c^{\dagger}(k,u)
  \label{eq:E4}
\end{equation}

$g'(k,u) = u c'(k,u) - v c(-k,d)$

\begin{equation}
  g^{\dagger}(k,u) = u c^{\dagger}(k,u) - v c(-k,d)
  \label{eq:E5}
\end{equation}

$g'(-k,d) = u c'(-k,d) + v c(k,u)$

\begin{equation}
  g^{\dagger}(-k,d) = u c^{\dagger}(-k,d) + v c(k,u)
  \label{eq:E6}
\end{equation}

4本を電子演算子について解き直し（逆変換）、準粒子演算子で電子演算子を表す。

E3, E4, E5, E6 を c(k,u), c(-k,d), c'(k,u), c'(-k,d) について逆変換して

\begin{equation}
  c_{k,u} = \frac{g^{\dagger}_{-k,d} \cdot v + g_{k,u} \cdot u}{u^{2} + v^{2}}
  \label{eq:E7}
\end{equation}

\begin{equation}
  c_{-k,d} = \frac{ - g^{\dagger}_{k,u} \cdot v + g_{-k,d} \cdot u}{u^{2} + v^{2}}
  \label{eq:E8}
\end{equation}

\begin{equation}
  c^{\dagger}_{k,u} = \frac{g^{\dagger}_{k,u} \cdot u + g_{-k,d} \cdot v}{u^{2} + v^{2}}
  \label{eq:E9}
\end{equation}

\begin{equation}
  c^{\dagger}_{-k,d} = \frac{g^{\dagger}_{-k,d} \cdot u - g_{k,u} \cdot v}{u^{2} + v^{2}}
  \label{eq:E10}
\end{equation}

逆変換をハミルトニアンへ代入する。対生成項の係数がゼロになるよう $u,v$ を選ぶのが Bogoliubov 条件。

E1 の c(k,u) に E7 を、c(-k,d) に E8 を代入する

\begin{equation}
  H_BCS = \left(2 \frac{v^{2} \xi}{\left( u^{2} + v^{2} \right)^{2}} + \frac{2 \Delta u v}{\left( u^{2} + v^{2} \right)^{2}}\right)  + \left( - \frac{u v \xi}{\left( u^{2} + v^{2} \right)^{2}} + \frac{ - u v \xi}{\left( u^{2} + v^{2} \right)^{2}} - \Delta \left( \frac{ - v^{2}}{\left( u^{2} + v^{2} \right)^{2}} + \frac{u^{2}}{\left( u^{2} + v^{2} \right)^{2}} \right)\right) {g}_{-k,d}^{\dagger} {g}_{ku}^{\dagger} + \left(\frac{u^{2} \xi}{\left( u^{2} + v^{2} \right)^{2}} - \frac{v^{2} \xi}{\left( u^{2} + v^{2} \right)^{2}} - \frac{2 \Delta u v}{\left( u^{2} + v^{2} \right)^{2}}\right) {g}_{-k,d}^{\dagger} {g}_{-k,d} + \left(\frac{u^{2} \xi}{\left( u^{2} + v^{2} \right)^{2}} + \frac{ - 2 \Delta u v}{\left( u^{2} + v^{2} \right)^{2}} + \frac{ - v^{2} \xi}{\left( u^{2} + v^{2} \right)^{2}}\right) {g}_{ku}^{\dagger} {g}_{ku} + \left(2 \frac{u v \xi}{\left( u^{2} + v^{2} \right)^{2}} + \Delta \left( \frac{ - v^{2}}{\left( u^{2} + v^{2} \right)^{2}} + \frac{u^{2}}{\left( u^{2} + v^{2} \right)^{2}} \right)\right) {g}_{-k,d} {g}_{ku}
  \label{eq:E11}
\end{equation}

ここからギャップ方程式へ向かう。クーパー対の異常平均 $\langle c_{-k\downarrow}c_{k\uparrow}\rangle$ を求めるため、まず対の積を式として登録する。

c(-k,d) c(k,u)

\begin{equation}
  c(-k,d) c(k,u)
  \label{eq:E12}
\end{equation}

準粒子基底へ移す。

E12 の c(k,u) に E7 を、c(-k,d) に E8 を代入する

\begin{equation}
  c(-k,d) c(k,u) = \frac{u v}{\left( u^{2} + v^{2} \right)^{2}}  + \frac{v^{2}}{\left( u^{2} + v^{2} \right)^{2}} {g}_{-k,d}^{\dagger} {g}_{ku}^{\dagger} + \frac{ - u v}{\left( u^{2} + v^{2} \right)^{2}} {g}_{-k,d}^{\dagger} {g}_{-k,d} + \frac{ - u v}{\left( u^{2} + v^{2} \right)^{2}} {g}_{ku}^{\dagger} {g}_{ku} + \frac{u^{2}}{\left( u^{2} + v^{2} \right)^{2}} {g}_{-k,d} {g}_{ku}
  \label{eq:E13}
\end{equation}

準粒子の真空（＝BCS基底状態）での期待値を取る。ここで演算子代数は終わり、以降は c 数の計算になる。

\begin{equation}
  c(-k,d) c(k,u) = \frac{u v}{\left( u^{2} + v^{2} \right)^{2}}
  \label{eq:E14}
\end{equation}

演算子を終えたので、残りのスカラー代数は Maxima の領分。セッション境界を越える（fork であって切り替えではない——このセッションはそのまま残り、渡し先は新しいファイルになる）。渡るのは右辺の値と、それを渡し先で establish し直すための宣言（符号込み）。左辺の演算子ラベルはスカラー CAS で意味を持たないので、値が「何の期待値か」は渡し先のメモに残る。

E14 を Maxima セッションに渡して

\end{document}